\section{AIME II 2020}
        \begin{enumerate}
        	\item (\textit{AIME II 2020}) Find the number of ordered pairs of positive integers $(m,n)$ such that ${m^2n = 20 ^{20}}$.


 
 \textbf{Answer: }231



\textbf{Solution}

In this problem, we want to find the number of ordered pairs $(m, n)$ such that $m^2n = 20^{20}$. Let $x = m^2$. Therefore, we want two numbers, $x$ and $n$, such that their product is $20^{20}$ and $x$ is a perfect square. Note that there is exactly one valid $n$ for a unique $x$, which is $\frac{20^{20}}{x}$. This reduces the problem to finding the number of unique perfect square factors of $20^{20}$.

 
$20^{20} = 2^{40} \cdot 5^{20} = \left(2^2\right)^{20}\cdot\left(5^2\right)^{10}.$ Therefore, the answer is $21 \cdot 11 = \boxed{231}.$

 





\textbf{Solution 2 (Official MAA)}

Because $20^{20}=2^{40}5^{20}$, if $m^2n = 20^{20}$, there must be nonnegative integers $a$, $b$, $c$, and $d$ such that$m = 2^a5^b$ and $n = 2^c5^d$. Then
 \[2a + c = 40\]and
 \[2b+d = 20\]The first equation has $21$ solutions corresponding to $a = 0,1,2,\dots,20$, and the second equation has $11$ solutions corresponding to $b = 0,1,2,\dots,10$. Therefore there are a total of $21\cdot11 = \boxed{231}$ ordered pairs $(m,n)$ such that $m^2n = 20^{20}$.

 \textbf{Purple Comet Math Meet April 2020} Notice, that this was the exact same problem (with different wording of course) as Purple Comet HS problem 3 and remembering the answer, put $\boxed{231}$.

 \href{https://purplecomet.org/views/data/2020HSSolutions.pdf}{https://purplecomet.org/views/data/2020HSSolutions.pdf}

 




	\item (\textit{AIME II 2020}) Let $P$ be a point chosen uniformly at random in the interior of the unit square with vertices at $(0,0), (1,0), (1,1)$, and $(0,1)$. The probability that the slope of the line determined by $P$ and the point $\left(\frac58, \frac38 \right)$ is greater than or equal to $\frac12$ can be written as $\frac{m}{n}$, where $m$ and $n$ are relatively prime positive integers. Find $m+n$.


 
 \textbf{Answer: }171



\textbf{Solution}

The areas bounded by the unit square and alternately bounded by the lines through $\left(\frac{5}{8},\frac{3}{8}\right)$ that are vertical or have a slope of $1/2$ show where $P$ can be placed to satisfy the condition. One of the areas is a trapezoid with bases $1/16$ and $3/8$ and height $5/8$. The other area is a trapezoid with bases $7/16$ and $5/8$ and height $3/8$. Then, 
 \[\frac{\frac{1}{16}+\frac{3}{8}}{2}\cdot\frac{5}{8}+\frac{\frac{7}{16}+\frac{5}{8}}{2}\cdot\frac{3}{8}=\frac{86}{256}=\frac{43}{128}\implies43+128=\boxed{171}\]~mn28407

 





\textbf{Solution 2 (Official MAA)}

The line through the fixed point $\left(\frac58,\frac38\right)$ with slope $\frac12$ has equation $y=\frac12 x + \frac1{16}$. The slope between $P$ and the fixed point exceeds $\frac12$ if $P$ falls within the shaded region in the diagram below consisting of two trapezoids with area
 \[\frac{\frac1{16}+\frac38}2\cdot\frac58 + \frac{\frac58+\frac7{16}}2\cdot\frac38 = \frac{43}{128}.\]Because the entire square has area $1,$ the required probability is $\frac{43}{128}$. The requested sum is $43+128 = 171$.
                    \begin{figure}[H]
                        \centering
                        \includegraphics[width=0.5\textwidth]{Solution/AIME_2020_II_2_Solution_Figure_1.png}
                    \end{figure}
                    

 


 




	\item (\textit{AIME II 2020}) The value of $x$ that satisfies $\log_{2^x} 3^{20} = \log_{2^{x+3}} 3^{2020}$ can be written as $\frac{m}{n}$, where $m$ and $n$ are relatively prime positive integers. Find $m+n$.


 
 \textbf{Answer: }103



\textbf{Solution}

Let $\log _{2^x}3^{20}=\log _{2^{x+3}}3^{2020}=n$. Based on the equation, we get $(2^x)^n=3^{20}$ and $(2^{x+3})^n=3^{2020}$. Expanding the second equation, we get $8^n\cdot2^{xn}=3^{2020}$. Substituting the first equation in, we get $8^n\cdot3^{20}=3^{2020}$, so $8^n=3^{2000}$. Taking the 100th root, we get $8^{\frac{n}{100}}=3^{20}$. Therefore, $(2^{\frac{3}{100}})^n=3^{20}$, and using the our first equation($2^{xn}=3^{20}$), we get $x=\frac{3}{100}$ and the answer is $\boxed{103}$.~rayfish

 \textbf{Easiest Solution} Recall the identity $\log_{a^n} b^{m} = \frac{m}{n}\log_{a} b$ (which is easily proven using exponents or change of base).Then this problem turns into 
 \[\frac{20}{x}\log_{2} 3 = \frac{2020}{x+3}\log_{2} 3\]Divide $\log_{2} 3$ from both sides. And we are left with $\frac{20}{x}=\frac{2020}{x+3}$.Solving this simple equation we get 
 \[x = \frac{3}{100} \Rightarrow \boxed{103}\]~mlgjeffdoge21

 





\textbf{Solution 2}

Because $\log_a{b^c}=c\log_a{b},$ we have that $20\log_{2^x} 3 = 2020\log_{2^{x+3}} 3,$ or $\log_{2^x} 3 = 101\log_{2^{x+3}} 3.$ Since $\log_a{b}=\dfrac{1}{\log_b{a}},$ $\log_{2^x} 3=\dfrac{1}{\log_{3} 2^x},$ and $101\log_{2^{x+3}} 3=101\dfrac{1}{\log_{3}2^{x+3}},$ thus resulting in $\log_{3}2^{x+3}=101\log_{3} 2^x,$ or $\log_{3}2^{x+3}=\log_{3} 2^{101x}.$ We remove the base 3 logarithm and the power of 2 to yield $x+3=101x,$ or $x=\dfrac{3}{100}.$

 Our answer is $\boxed{3+100=103}.$~ OreoChocolate

 





\textbf{Solution 3 (Official MAA)}

Using the Change of Base Formula to convert the logarithms in the given equation to base $2$ yields
 \[\frac{\log_2 3^{20}}{\log_2 2^x} = \frac{\log_2 3^{2020}}{\log_2 2^{x+3}}, \text{~ and then ~} \frac{20\log_2 3}{x\cdot\log_2 2} = \frac{2020\log_2 3}{(x+3)\log_2 2}.\]Canceling the logarithm factors then yields
 \[\frac{20}x = \frac{2020}{x+3},\]which has solution $x = \frac3{100}.$ The requested sum is $3 + 100 = 103$.

 





\textbf{Solution 4}

$\log_{2^x} 3^{20} = 2^{xy} = 3^{20}$

 $\log_{2^{x+3}} 3^{2020} = (2^{x+3})^y = 3^{2020}$

 $3^{2020} = (3^{20})^{101}$

 $(2^{xy})^{101} = (2^{x+3})^y = 3^{2020}$

 $(2^{xy})^{101} = (2^{x+3})^y \Rightarrow 2^{101xy} = 2^{xy+3y} \Rightarrow 101xy = xy + 3y \Rightarrow 101xy = y(x+3)$

 $101x = x + 3$

 $100x = 3$

 $x = \frac{3}{100}$

 $100 + 3 = \boxed{103}$ ~Airplanes2007

 





\textbf{Solution 5}

$\textbf{Lemma}:$

 $\log_{a^c}{b} = \frac{1}{c} \log_{a}{b}$

 $\textbf{Proof of the Lemma}:$

 Change of base and exponent property of logarithms on LHS gives $\frac{\log{b}}{c\log{a}}$ which simplifies to $\frac{1}{c} \cdot \frac{\log{b}}{\log{a}}$ in which converse of change of base gives $\frac{1}{c} \log_{a}{b}$ as desired.

 $\textbf{2020 AIME II Q3}$

 Apply the lemma on both sides to get $\frac{1}{x} \cdot \log_{2}{3^{20}} = \frac{1}{x+3} \cdot \log_{2}{3^{2020}}$. Then, apply the exponent property to obtain $\log_{2}{3^{\frac{20}{x}}} = \log_{2}{3^{\frac{2020}{x+3}}}$. Equivalent logarithms give $3^{\frac{20}{x}} = 3^{\frac{2020}{x+3}}$ and equivalent bases give $\frac{20}{x} = \frac{2020}{x+3}~:~x \neq -3, 0$. Then, the statement $60+20x = 2020x$ is obtained following that $x = \frac{3}{100} \implies \boxed{103}$.

 




	\item (\textit{AIME II 2020}) Triangles $\triangle ABC$ and $\triangle A'B'C'$ lie in the coordinate plane with vertices $A(0,0)$, $B(0,12)$, $C(16,0)$, $A'(24,18)$, $B'(36,18)$, $C'(24,2)$. A rotation of $m$ degrees clockwise around the point $(x,y)$ where $0<m<180$, will transform $\triangle ABC$ to $\triangle A'B'C'$. Find $m+x+y$.


 
 \textbf{Answer: }108



\textbf{Solution}

After sketching, it is clear a $90^{\circ}$ rotation is done about $(x,y)$. Looking between $A$ and $A'$, $x+y=18$. Thus $90+18=\boxed{108}$.~mn28407

 





\textbf{Solution 2 (Official MAA)}

Because the rotation sends the vertical segment $\overline{AB}$ to the horizontal segment $\overline{A'B'}$, the angle of rotation is $90^\circ$ degrees clockwise. For any point $(x,y)$ not at the origin, the line segments from $(0,0)$ to $(x,y)$ and from $(x,y)$ to $(x-y,y+x)$ are perpendicular and are the same length. Thus a $90^\circ$ clockwise rotation around the point $(x,y)$ sends the point $A(0,0)$ to the point $(x-y,y+x) = A'(24,18)$. This has the solution $(x,y) = (21,-3)$. The requested sum is $90+21-3=108$.

 





\textbf{Solution 3}


                    \begin{figure}[H]
                        \centering
                        \includegraphics[width=0.5\textwidth]{Solution/AIME_2020_II_4_Solution_Figure_1.png}
                    \end{figure}
                    

 We first draw a diagram with the correct Cartesian coordinates and a center of rotation $P$. Note that $PC=PC'$ because $P$ lies on the perpendicular bisector of $CC'$ (it must be equidistant from $C$ and $C'$ by properties of a rotation). 

 Since $AB$ is vertical while $A'B'$ is horizontal, we have that the angle of rotation must be $90^{\circ}$, and therefore $\angle P = 90^{\circ}$. Therefore, $CPC'$ is a 45-45-90 right triangle, and $CD=DP$. 

 We calculate $D$ to be $(20,1)$. Since we translate $4$ right and $1$ up to get from point $C$ to point $D$, we must translate $1$ right and $4$ down to get to point $P$. This gives us $P(21,-3)$. Our answer is then $90+21-3=\boxed{108}$. ~Lopkiloinm & samrocksnature

 





\textbf{Solution 4}

For the above reasons, the transformation is simply a $90^\circ$ rotation. Proceed with complex numbers on the points $C$ and $C'$. Let $(x, y)$ be the origin. Thus, $C \rightarrow (16-x)+(-y)i$ and $C' \rightarrow (24-x)+(2-y)i$. The transformation from $C'$ to $C$ is a multiplication of $i$, which yields $(16-x)+(-y)i=(y-2)+(24-x)i$. Equating the real and complex terms results in the equations $16-x=y-2$ and $-y=24-x$. Solving, $(x, y) : (21, -3) \rightarrow 90+21-3=\boxed{108}$

 





\textbf{Solution 5}

We know that the rotation point $P$ has to be equidistant from both $A$ and $A'$ so it has to lie on the line that is on the midpoint of the segment $AA'$ and also the line has to be perpendicular to $AA'$. Solving, we get the line is $y=\frac{-4}{3}x+25$. Doing the same for $B$ and $B'$, we get that $y=-6x+123$. Since the point $P$ of rotation must lie on both of these lines, we set them equal, solve and get: $x=21$,$y=-3$. We can also easily see that the degree of rotation is $90$ since $AB$ is initially vertical, and now it is horizontal. Also, we can just sketch this on a coordinate plane and easily realize the same. Hence, the answer is $21-3+90 = \boxed{108}$

 


 





\textbf{Solution 6}


                    \begin{figure}[H]
                        \centering
                        \includegraphics[width=0.5\textwidth]{Solution/AIME_2020_II_4_Solution_Figure_2.png}
                    \end{figure}
                    We make transformation of line $AB$ into line $A'B'$ using axes symmetry. Point $D(0,18)$ is the crosspoint of this lines. Equation of line $DO$ is 
 \[x + y = 18.\]$\triangle ABC$ maps into $\triangle A''B''C''$ where 
 \[A''(18,18), B''(6,18), C''(18,2).\]

 We make transform of the line $A''C''$ into line $A'C'$ using axes symmetry with respect to line
 \[x  = \frac {A'' + A'}{2} = \frac {24 + 18}{2} = 21.\]The composition of two axial symmetries is a rotation through an angle twice as large as the angle between the axes $(45^o)$ around the point of their intersection $O(21, – 3).$ 
 \[m + x + y = 90 + 21 – 3 = \boxed {108}\].

 \textbf{vladimir.shelomovskii@gmail.com, vvsss}

 





\textbf{Solution 7 (Matrix and Transformations)}

For a matrix to rotate a figure on a coordinate plane by $m$ degrees, it is written as:$\left[ {\begin{array}{cc}  \cos(m^{\circ}) & \sin(m^{\circ})  \\  -\sin(m^{\circ}) & \cos(m^{\circ})  \\  \end{array} } \right]$

 We can translate the whole figure so that the centre of rotation is at $(0,0)$, which is equivalent to subtracting $x$ and $y$ from all $x$-coordinates and the $y$-coordinates respectively of the given points.

 We then record all the points $A$, $B$, $C$ in a matrix as follows:$\left[ {\begin{array}{ccc}  0-x & 0-x & 16-x  \\  0-y & 12-y & 0-y  \\  \end{array} } \right]$

 and all the points $A'$, $B'$, $C'$ in a matrix as follows:$\left[ {\begin{array}{ccc}  24-x & 36-x & 24-x  \\  18-y & 18-y & 2-y  \\  \end{array} } \right]$

 Since $\triangle A'B'C'$ is a rotation around $(x,y)$ of $\triangle ABC$ by $m^{\circ}$, by the left multiplication rule, we can equate that:

 $\left[ {\begin{array}{cc}  \cos(m^{\circ}) & \sin(m^{\circ})  \\  -\sin(m^{\circ}) & \cos(m^{\circ})  \\  \end{array} } \right]$$\left[ {\begin{array}{ccc}  0-x & 0-x & 16-x  \\  0-y & 12-y & 0-y  \\  \end{array} } \right]$$=$$\left[ {\begin{array}{ccc}  24-x & 36-x & 24-x  \\  18-y & 18-y & 2-y  \\  \end{array} } \right]$

 We can obtain the follow equations:$\begin{cases} -x\cos(m^{\circ})-y\sin(m^{\circ})=24-x \\ -x\cos(m^{\circ})-y\sin(m^{\circ})+12\sin(m^{\circ})=36-x \\ x\sin(m^{\circ})-y\cos(m^{\circ})=24-x \end{cases}$

 From the first 2 equations, we get $m=90$, substituting into the 3rd equation, we get $x+y=18$.

 Therefore, $m+x+y=90+18=\boxed{108}$

 





\textbf{Solution 8}

It is clear that $\bigtriangleup CPC'$ is a $45-45-90$ right triangle so $m=90$. We use the $\tan$ angle formula, $\tan{(a-b)}=\frac{\tan(a)-\tan(b)}{1+\tan(a)\tan(b)}$ to find the slope of line $CP$. We know that line $CC'$ has slope $\frac{1}{4}$ and let $b=-45^{\circ}$, then plugging both values into the formula, we find that the slope of $CP$ is $\frac{-3}{5}$. Also, $CC'$ has length $\sqrt{34}$. Create a right triangle $KCP$ where $KP$ is parallel to the $x$ axis and $CP$ is the hypotenuse. Then $CK=3x$ and $KP=5x$ and doing Pythagorean on $\bigtriangleup KCP$ gives $x=1$. Therefore, we know that $P$ is a translation 3 units down and 5 units right from $C(16,0)$, from which we obtain $P(21,-3)$. Adding the three variables, we obtain $90+21-3=\boxed{108}$

 




	\item (\textit{AIME II 2020}) For each positive integer $n$, let $f(n)$ be the sum of the digits in the base-four representation of $n$ and let $g(n)$ be the sum of the digits in the base-eight representation of $f(n)$. For example, $f(2020) = f(133210_{\text{4}}) = 10 = 12_{\text{8}}$, and $g(2020) = \text{the digit sum of }12_{\text{8}} = 3$. Let $N$ be the least value of $n$ such that the base-sixteen representation of $g(n)$ cannot be expressed using only the digits $0$ through $9$. Find the remainder when $N$ is divided by $1000$.


 
 \textbf{Answer: }151



\textbf{Solution 1}

Let's work backwards. The minimum base-sixteen representation of $g(n)$ that cannot be expressed using only the digits $0$ through $9$ is $A_{16}$, which is equal to $10$ in base 10.  Thus, the sum of the digits of the base-eight representation of the sum of the digits of $f(n)$ is $10$.  The minimum value for which this is achieved is $37_8$.  We have that $37_8 = 31$.  Thus, the sum of the digits of the base-four representation of $n$ is $31$.  The minimum value for which this is achieved is $13,333,333,333_4$.  We just need this value in base 10 modulo 1000.  We get $13,333,333,333_4 = 3(1 + 4 + 4^2 + \dots + 4^8 + 4^9) + 4^{10} = 3\left(\dfrac{4^{10} - 1}{3}\right) + 4^{10} = 2*4^{10} - 1$.  Taking this value modulo $1000$, we get the final answer of $\boxed{151}$.  (If you are having trouble with this step, note that $2^{10} = 1024 \equiv 24 \pmod{1000}$) ~ TopNotchMath

 





\textbf{Solution 2 (Official MAA)}

First note that if $h_b(s)$ is the least positive integer whose digit sum, in some fixed base $b$, is $s$, then $h_b$ is a strictly increasing function. This together with the fact that $g(N) \ge 10$ shows that $f(N)$ is the least positive integer whose base-eight digit sum is 10. Thus $f(N) = 37_\text{eight} = 31$, and $N$ is the least positive integer whose base-four digit sum is $31.$ Therefore 
 \[N = 13333333333_\text{four} = 2\cdot4^{10} - 1 = 2\cdot1024^2 - 1\]
 \[\equiv 2\cdot24^2 - 1 \equiv 151 \pmod{1000}.\]

 




	\item (\textit{AIME II 2020}) Define a sequence recursively by $t_1 = 20$, $t_2 = 21$, and
 \[t_n = \frac{5t_{n-1}+1}{25t_{n-2}}\]for all $n \ge 3$. Then $t_{2020}$ can be written as $\frac{p}{q}$, where $p$ and $q$ are relatively prime positive integers. Find $p+q$.


 
 \textbf{Answer: }626



\textbf{Solution}

Let $t_n=\frac{s_n}{5}$. Then, we have $s_n=\frac{s_{n-1}+1}{s_{n-2}}$ where $s_1 = 100$ and $s_2 = 105$. By substitution, we find $s_3 = \frac{53}{50}$, $s_4=\frac{103}{105\cdot50}$, $s_5=\frac{101}{105}$, $s_6=100$, and $s_7=105$. So $s_n$ has a period of $5$. Thus $s_{2020}=s_5=\frac{101}{105}$. So, $\frac{101}{105\cdot 5}\implies 101+525=\boxed{626}$.~mn28407

 





\textbf{Solution 2 (Official MAA)}

More generally, let the first two terms be $a$ and $b$ and replace $5$ and $25$ in the recursive formula by $k$ and $k^2$, respectively. Then some algebraic calculation shows that
 \[t_3 = \frac{b\,k+1}{a\, k^2},~~t_4 = \frac{a\, k + b\,k+1}{a\,b\, k^3},~~ t_5 = \frac{a\,k+1}{b\, k^2},~~ t_6 = a, \text{~ and ~}t_7 =b,\]so the sequence is periodic with period $5$. Therefore
 \[t_{2020} = t_{5} =  \frac{20\cdot 5 +1}{21\cdot  25} = \frac{101}{525}.\]The requested sum is $101+525=626$.

 





\textbf{Solution 3}

Let us examine the first few terms of this sequence and see if we can find a pattern. We are obviously given $t_1 = 20$ and $t_2 = 21$, so now we are able to determine the numerical value of $t_3$ using this information:
 \[t_3 = \frac{5t_{3-1}+1}{25t_{3-2}} = \frac{5t_{2}+1}{25t_{1}} = \frac{5(21) + 1}{25(20)} = \frac{105 + 1}{500}t_3 = \frac{106}{500}  = \frac{53}{250}\]
 \[t_4 = \frac{5t_{4-1}+1}{25t_{4-2}} = \frac{5t_{3}+1}{25t_{2}} = \frac{5(\frac{53}{250}) + 1}{25(21)} = \frac{\frac{53}{50} + 1}{525} = \frac{\frac{103}{50}}{525} = \frac{103}{26250}\]
 \[t_5 = \frac{5t_{5-1}+1}{25t_{5-2}} = \frac{5t_{4}+1}{25t_{3}} = \frac{5(\frac{103}{26250}) + 1}{25(\frac{53}{250})} = \frac{\frac{103}{5250} + 1}{\frac{53}{10}} = \frac{\frac{5353}{5250}}{\frac{53}{10}} = \frac{101}{525}\]
 \[t_6 = \frac{5t_{6-1}+1}{25t_{6-2}} = \frac{5t_{5}+1}{25t_{4}} = \frac{5(\frac{101}{525}) + 1}{25(\frac{103}{26250})} = \frac{\frac{101}{105} + 1}{\frac{103}{1050}} = \frac{\frac{206}{105}}{\frac{103}{1050}} \implies t_6 = 20\]

 Alas, we have figured this sequence is period 5! But since $2020 \equiv 5 \pmod 5$, we can state that $t_{2020} = t_5 = \frac{101}{525}$. According to the original problem statement, our answer is $\boxed{626}$. ~ nikenissan

 





\textbf{Solution 4}

Like solution 1, we will make a substitution. Let $u_n = 5t_n$. Then

 
 \[u_n = \frac{25t_{n-1} + 5}{25t_{n-2}} = \frac{u_{n-1} + 1}{u_{n-2}}\].

 Let $u_1 = a$, $u_2 = b$. The following holds for all $n \ge 3$. Then,

 
 \[u_3 = \frac{b+1}{a}\]
 \[u_4 = \frac{\frac{b+1+a}{a}}{b} = \frac{b+1+a}{ab}\]
 \[u_5 = \frac{\frac{b+1+a+ab}{ab}}{\frac{b+1}{a}} = \frac{\frac{(b+1)(a+1)}{ab}}{\frac{b+1}{a}} = \frac{(b+1)(a+1)}{ab} \cdot \frac{a}{b+1} = \frac{a+1}{b}\]
 \[u_6 = \frac{\frac{a+1+b}{b}}{\frac{b+1+a}{ab}} = \frac{a+1+b}{b} \cdot \frac{ab}{b+1+a} = a\].

 Thus, $u_n$ is periodic with periodicity 5.

 Note that $u_{2020} = u_5$, and thus $t_{2020} = t_5$. Also, $u_1 = 5t_1 = 100 = a$ and $u_2 = 5t_2 = 105 = b$. Then,

 
 \[5t_5 = \frac{101}{105}\]
 \[t_5 = \frac{101}{525}\]

 And thus $t_{2020} = t_5 = \frac{101}{525}$. The answer is $101 + 525 = \boxed{626}$.

 \textbf{Quick way to notice recursion loop} Round the first two values to both be 20. Then, the next element can be rounded to $\frac{1}{5}$. $t_4$ can then be quickly calculated to around $\frac{1}{250}$, and $t_5$ can be rounded to $\frac{1}{5}$. $t_6$ turns out to be around 20, which means that there is probably a loop with period 5. The rest of the solution proceeds normally.

 




	\item (\textit{AIME II 2020}) Two congruent right circular cones each with base radius $3$ and height $8$ have axes of symmetry that intersect at right angles at a point in the interior of the cones a distance $3$ from the base of each cone. A sphere with radius $r$ lies within both cones. The maximum possible value of $r^2$ is $\frac{m}{n}$, where $m$ and $n$ are relatively prime positive integers. Find $m+n$.


 
 \textbf{Answer: }298



\textbf{Solution (Official MAA)}

Consider the cross section of the cones and sphere by a plane that contains the two axes of symmetry of the cones as shown below. The sphere with maximum radius will be tangent to the sides of each of the cones. The center of that sphere must be on the axis of symmetry of each of the cones and thus must be at the intersection of their axes of symmetry. Let $A$ be the point in the cross section where the bases of the cones meet, and let $C$ be the center of the sphere. Let the axis of symmetry of one of the cones extend from its vertex, $B$, to the center of its base, $D$. Let the sphere be tangent to $\overline{AB}$ at $E$. The right triangles $\triangle ABD$ and $\triangle CBE$ are similar, implying that the radius of the sphere is
 \[CE = AD \cdot\frac{BC}{AB} = AD \cdot\frac{BD-CD}{AB} =3\cdot\frac5{\sqrt{8^2+3^2}} = \frac{15}{\sqrt{73}}=\sqrt{\frac{225}{73}}.\]The requested sum is $225+73=298$.
                    \begin{figure}[H]
                        \centering
                        \includegraphics[width=0.5\textwidth]{Solution/AIME_2020_II_7_Solution_Figure_1.png}
                    \end{figure}
                    

 Not part of MAA's solution, but this: \href{https://www.geogebra.org/calculator/xv4nm97a}{https://www.geogebra.org/calculator/xv4nm97a} is a good visual of the cones in GeoGebra.

 





\textbf{Solution (Incircles + Coordinate Bash)}


                    \begin{figure}[H]
                        \centering
                        \includegraphics[width=0.5\textwidth]{Solution/AIME_2020_II_7_Solution_Figure_2.png}
                    \end{figure}
                    

 Let $A$ be the origin in the above diagram. Then $B$ is $(6,0)$, $C$ is $(0,6)$, $D$ is $(3,8)$, and $E$ is $(8,3)$.  Also, it is easy to see that the inscribed sphere is simply the inscribed circle of $AEF$. Then we want to find the intersection of $AD$ and $CE$ to determine the coordinates of point $F$. Note that $AD$ is line $y=\frac{8}{3}x$, and $CE$ is line $y=-\frac{3}{8}x+6$. Then, you can see that these lines are perpendicular, indicating that $AEF$ is a right triangle with right angle at $F$. Finding the intersection of the lines by solving the system, we get that $F$ is the point $(\frac{144}{73},\frac{384}{73})$ in this plane. Then, we can find the distances $EF$ and $AF$ by the distance formula, which are $\frac{55}{\sqrt{73}}$ and $\frac{48}{\sqrt{73}}$ respectively. Also, $AE=\sqrt{73}$. Then, to find the radius of this triangle's incircle, we use the formula $a=rs$ from which we get that $r=\frac{15}{\sqrt{73}}$ and $r^2=\frac{225}{73} \implies \boxed{298}$ is the answer.

 





\textbf{Solution (Clean analytic geometry)}

Using the diagram above, we notice that the desired length is simply the distance between the point $C$ and $\overline{AB}$. We can mark $C$ as $(3,3)$ since it is $3$ units away from each of the bases. Point $B$ is $(8,3)$. Thus, line $\overline{AB}$ is $y = \frac{3}{8}x \Rightarrow 3x + 8y = 0$. We can use the distance from point to line formula $\frac{|Ax_0 + By_0 + C|}{\sqrt{A^2 + B^2}}$, where $x_0$ and $y_0$ are the coordinates of the point, and A, B, and C are the coefficients of the line in form $Ax + By + C = 0$. Plugging everything in, we get 
 \[\frac{|3(3) - 8(3)|}{\sqrt{8^2 + 3^2}} = \frac{15}{\sqrt{73}} \Rightarrow \frac{225}{73} \Rightarrow \boxed{298}\].

 





\textbf{Solution 1: Graph paper coordbash}

We graph this on graph paper, with the scale of $\sqrt{2}:1$. So, we can find $OT$ then divide by $\sqrt{2}$ to convert to our desired units, then square the result. With 5 minutes' worth of coordbashing, we finally arrive at $298$.

 
                    \begin{figure}[H]
                        \centering
                        \includegraphics[width=0.5\textwidth]{Solution/AIME_2020_II_7_Solution_Figure_3.png}
                    \end{figure}
                    

 





\textbf{Solution 2}

Using the graph drawn above in MAA solution, we see that $\sqrt{AC^2-r^2}$ + $\sqrt{BC^2-r^2}$ = $AB$. $AB$ = $\sqrt{8^2+3^2}$ = $\sqrt{73}$, $AC$ = 3$\sqrt{2}$, $AB$ = 5. Plugging in and solve for $r^2$ gives us $r^2$ = $\frac{225}{73}$.

 




	\item (\textit{AIME II 2020}) Define a sequence recursively by $f_1(x)=|x-1|$ and $f_n(x)=f_{n-1}(|x-n|)$ for integers $n>1$. Find the least value of $n$ such that the sum of the zeros of $f_n$ exceeds $500,000$.


 
 \textbf{Answer: }101



\textbf{Solution 1 (Official MAA)}

First it will be shown by induction that the zeros of $f_n$ are the integers$a, {a+2,} {a+4,} \dots, {a + n(n-1)}$, where $a = n - \frac{n(n-1)}2.$

 This is certainly true for $n=1$. Suppose that it is true for $n = m-1 \ge 1$, and note that the zeros of $f_m$ are the solutions of $|x - m| = k$, where $k$ is a nonnegative zero of $f_{m-1}$. Because the zeros of $f_{m-1}$ form an arithmetic sequence with common difference $2,$ so do the zeros of $f_m$. The greatest zero of $f_{m-1}$ is
 \[m-1+\frac{(m-1)(m-2)}2 =\frac{m(m-1)}2,\]so the greatest zero of $f_m$ is $m+\frac{m(m-1)}2$ and the least is $m-\frac{m(m-1)}2$.

 It follows that the number of zeros of $f_n$ is $\frac{n(n-1)}2+1=\frac{n^2-n+2}2$, and their average value is $n$. The sum of the zeros of $f_n$ is
 \[\frac{n^3-n^2+2n}2.\]Let $S(n)=n^3-n^2+2n$, so the sum of the zeros exceeds $500{,}000$ if and only if $S(n) > 1{,}000{,}000 = 100^3\!.$ Because $S(n)$ is increasing for $n > 2$, the values $S(100) = 1{,}000{,}000 - 10{,}000 + 200 = 990{,}200$ and $S(101)=1{,}030{,}301 - 10{,}201 + 202 = 1{,}020{,}302$ show that the requested value of $n$ is $\boxed{101}$.

 





\textbf{Solution 2 (Same idea, easier to see)}

Starting from $f_1(x)=|x-1|$, we can track the solutions, the number of solutions, and their sum: 

 
 \[\begin{array}{c|c|c|c} n&Solutions&number&sum\\  1&1&1&1\\   2&1,3&2&4\\    3&0,2,4,6&4&12\\ 4&-2,0,2...10&7&28\\ 5&-5,-3,-1...15&11&55\\ \end{array}\]

 It is clear that the solutions form arithmetic sequences with a difference of 2, and the negative solutions cancel out all but $n$ of the $1+\frac{n(n-1)}{2}$ solutions. Thus, the sum of the solutions is $n \cdot [1+\frac{n(n-1)}{2}]$, which is a cubic function.(Side Note: Gergor-Newton Interpolation Formula is applicable here)

 $n \cdot [1+\frac{n(n-1)}{2}]>500,000$

 Multiplying both sides by $2$, 

 $n \cdot [2+n(n-1)]>1,000,000$

 1 million is $10^6=100^3$, so the solution should be close to $100$.

 100 is slightly too small, so $\boxed{101}$ works. 

 




	\item (\textit{AIME II 2020}) While watching a show, Ayako, Billy, Carlos, Dahlia, Ehuang, and Frank sat in that order in a row of six chairs. During the break, they went to the kitchen for a snack. When they came back, they sat on those six chairs in such a way that if two of them sat next to each other before the break, then they did not sit next to each other after the break. Find the number of possible seating orders they could have chosen after the break.


 
 \textbf{Answer: }090



\textbf{Solution 1 (Official MAA)}

Ayako ($A$), Billy $(B)$, Carlos $(C)$, Dahlia $(D)$, Ehuang $(E)$, and Frank $(F)$ originally sat in the order $ABCDEF$.Let $T(XY)$ denote the set of seatings where $X$ and $Y$ sit next to each other after the break. Then the required number of seating orders is given by the Inclusion-Exclusion Principle as
 \[6!-\big(|T(AB)|+|T(BC)|+|T(CD)|+|T(DE)|+|T(EF)|\big)+\]
 \[\big(|T(AB)\cap T(BC)|+|T(AB)\cap T(CD)|+\cdots\big) - \cdots.\]Each term can be calculated separately.

 (a) $|T(AB)|=|T(BC)|=|T(CD)|=|T(DE)|=|T(EF)|=2\cdot5!=240.$ Because there are $5$ terms, the sum is $5\cdot240=1200$.

 (b) For $|T(XY)\cap T(ZW)|$, if $Y=Z$, then $XYW$ must sit consecutively, so $|T(XY)\cap T(ZW)|=2\cdot4!=48$. There are $4$ terms that satisfy $Y=Z$, so the sum is $4\cdot 48=192$. If $XY$ and $ZW$ are pairwise disjoint, then $|T(XY)\cap T(ZW)|=2^2\cdot4!=96$. There are $6$ terms, so the sum is $6\cdot96=576$.

 (c) If there are at least three pairs that sit next to each other, consider these three subcases:If the three pairs are consecutive, the sum is $3\cdot2\cdot3!=36$.If exactly two of the pairs are consecutive, the sum is $6\cdot2^2\cdot3!=144$.If none of the three pairs is consecutive, the sum is $1\cdot2^3\cdot3!=48.$

 (d) If there are at least four pairs that sit next to each other, then if the pairs are consecutive, the sum is $2\cdot2\cdot2!=8$. If the pairs are not consecutive, then the sum is $3\cdot2^2\cdot2!=24$.

 (e) If all five pairs sit next to each other, the number is $1\cdot2\cdot1!=2$.

 Therefore the required number of seating orders is
 \[6!-1200+(192+576)-{(36+144+48)+(8+24)-2}=90.\]

 Note: See the \href{https://oeis.org/A002464}{A002464 of the On-Line Encyclopedia of Integer Sequences} for equivalent formulations.

 





\textbf{Solution 2 (Casework, fast method)}

We proceed with casework based on the person who sits first after the break.

 $\textbf{Case 1:}$ A is first. Then the first three people in the row can be ACE, ACF, ADB, ADF, AEB, AEC, AFB, AFC, or AFD, which yield 2, 1, 2, 2, 1, 2, 0, 1, and 1 possible configurations, respectively, implying 2 + 1 + 2 + 2 + 1 + 2 + 0 + 1 + 1 = 12 possible configurations in this case.

 $\textbf{Case 2:}$ B is first. Then the first three people in the row can be BDA, BDF, BEA, BEC, BFA, BFC, or BFD, which yield 2, 4, 2, 4, 0, 1, and 2 possible configurations, respectively, implying 2 + 4 + 2 + 4 + 0 + 1 + 2 = 15 possible configurations in this case.

 $\textbf{Case 3:}$ C is first. Then the first three people in the row can be CAD, CAE, CAF, CEA, CEB, CFA, CFB, or CFD, which yield 1, 2, 1, 4, 4, 2, 2, and 2 possible configurations, respectively, implying 1 + 2 + 1 + 4 + 4 + 2 + 2 + 2 = 18 possible configurations in this case.

 Finally, the number of valid configurations for A and F, B and E, and C and D are equal by symmetry, so our final answer is 2(12 + 15 + 18), which computes to be $\boxed{090}.$ ~peace09

 





\textbf{Solution 3}

We determine the order of A, B, C, relative to each other. Then, we will insert D, E, F into the alignment and calculate the total number of possibilities. This solution can be visualised as standing in lines rather than sitting on chairs.

 There are 6 possible alignments for A, B, and C, but we only evaluate 3 because the other 3 cases can be mirrored to overlap these 3 cases.

 $\textbf{Case 1: A B C}$ In this case, there must be a person standing between A and B and also between B and C. Also, D cannot be adjacent to C. There are 9 possibilities.

 $\textbf{Case 2: A C B}$ In this case, there must be a person standing between B and C. Also, D cannot be adjacent to C. There are 12 possibilities.

 $\textbf{Case 3: C A B}$ In this case, there must be a person standing between A and B. Also, D cannot be adjacent to C. There are 24 possibilities.

 So the total number of cases is 2(9+12+24)=$\boxed{090}$.

 -Superdolphin~xHypotenuse

 





\textbf{Solution 4 (Recursion)}

Consider arrangements of numbers $1, 2, 3, ..., n$.

 Let $f(n,k)$ be the number of arrangements in which $1$ and $2$, $2$ and $3$, $3$ and $4$, ..., $k-1$ and $k$ aren't together. We need to find $f(6,6)$.

 Let $d(n,k)$ be the number of arrangements in which $1$ and $2$, $2$ and $3$, $3$ and $4$, ..., $k-2$ and $k-1$ aren't together, but $(k-1)$ and $k$ are together.

 Then, 

 
 \[f(n,k+1) = f(n,k) - d(n,k+1)  ......(1)\]

 
 \[d(n,k+1) = 2f(n-1,k) + d(n-1,k)  ......(2)\]

 Hence,
 \[f(n,k+1)=f(n,k)-f(n-1,k)-f(n-1,k-1)\]

 And because $f(n,0) = f(n,1) = n!$, it's easy to get $f(6,6) = \boxed{090}$

 

 \[\begin{tabular}{|c|c|c|c|c|c|c|}\hline k,n & n=1 & n=2 & n=3 & n=4 & n=5 & n=6\\\hline k=1 & 1 & 2 & 6 & 24 & 120 & 720\\\hline k=2 & & 0 & 2 & 12 & 72 & 480\\\hline k=3 & & & 0 & 4 & 36 & 288\\\hline k=4 & & & & 2 & 20 & 180\\\hline k=5 & & & & & 14 & 124\\\hline k=6 & & & & & & 90\\\hline \end{tabular}\]

 
Note: How to get the two equations (1) and (2).

 
1. If $1$ and $2$, $2$ and $3$, ... , $k-1$ and $k$ aren't together, there are two cases:

 (a) $k$ and $k+1$ aren't together. There are $f(n,k+1)$ ways of arrangements.

 (b) $k$ and $k+1$ are together. There are $d(n,k+1)$ ways of arrangements.

 Thus, $f(n,k)=f(n,k+1)+d(n,k+1)$, i.e. $f(n,k+1)=f(n,k)-d(n,k+1)$, which is the first equation.

 
2. If $1$ and $2$, $2$ and $3$, ... , $k-1$ and $k$ aren't together, but $k$ and $k+1$ are together, there are 2 cases : 

 (a) $k-1$ isn't together with $k$ or $k+1$. We first determine the order of $k$ and $k+1$, then, we can treat them as one number since they should be always together. Hence, the number of arrangements is $2f(n-1,k)$.

 (b) $k-1$ is together with $k$ or $k+1$. We first determine the order of $k$ and $k+1$, then  treat them as one "super number". Since $k-1$ is together with the "super number", there are $2d(n-1,k)$ ways of arrangements. However, $k-1$ can only be on one side of the "super number" because $k-1$ and $k$ aren't together, so we need to multiple a $1/2$ to the number of arrangements. Finally, there are $2d(n-1,k)*(1/2)=d(n-1,k)$ ways of arrangements.

 Thus, $d(n,k+1)=2f(n-1,k)+d(n-1,k)$, which is the second equation.

 


 


 


 




	\item (\textit{AIME II 2020}) Find the sum of all positive integers $n$ such that when $1^3+2^3+3^3+\cdots +n^3$ is divided by $n+5$, the remainder is $17$.


 
 \textbf{Answer: }239



\textbf{Solution 1}

The formula for the sum of cubes, also known as Nicomachus's Theorem, is as follows:
 \[1^3+2^3+3^3+\dots+k^3=(1+2+3+\dots+k)^2=\left(\frac{k(k+1)}{2}\right)^2\]for any positive integer $k$. 

 So let's apply this to this problem.

 Let $m=n+5$. Then we have
 \begin{align*} 1^3+2^3+3^3+\dots+(m-5)^3&\equiv 17 \mod m \\ \left(\frac{(m-5)(m-4)}{2}\right)^2&\equiv 17 \mod m \\ \left(\dfrac{m(m-9)+20}2\right)^2&\equiv 17\mod m \\ \left(m(m-9)+20\right)^2&\equiv 4\cdot 17\mod m \\ \left(20\right)^2&\equiv 68\mod m \\ 332 &\equiv 0 \mod m \\ \end{align*}
So, $m\in\{83,166,332\}$. Testing the cases, only $332$ fails. This leaves $78+161=\boxed{239}$. 

 $\LaTeX$ and formatting adjustments and intermediate steps for clarification by Technodoggo.

 





\textbf{Solution 2 (Official MAA 1)}

The sum of the cubes from 1 to $n$ is
 \[1^3+2^3+\cdots+n^3=\frac{n^2(n+1)^2}{4}.\]For this to be equal to $(n+5)q+17$ for some integer $q$, it must be that
 \[n^2(n+1)^2=4(n+5)q+4\cdot 17,\]so
 \[n^2(n+1)^2 \equiv 4 \cdot 17= 68\hskip-.2cm \pmod{n+5}.\]But $n^2(n+1)^2 \equiv (-5)^2(-4)^2 = 400 \pmod{n+5}.$ Thus $n^2(n+1)^2$ is congruent to both $68$ and $400,$ which implies that $n+5$ divides $400-68 = 332=2^2 \cdot 83$. Because $n+5 > 17$, the only choices for $n+5$ are $83, 166,$ and $332.$ Checking all three cases verifies that $n=78$ and $n=161$ work, but $n=327$ does not. The requested sum is $78+161 = 239$.

 





\textbf{Solution 3 (Official MAA 2)}

The sum of the cubes of the integers from $1$ through $n$ is
 \[1^3+2^3+\cdots+n^3=\frac{n^2(n+1)^2}{4},\]which, when divided by $n+5$, has quotient
 \[Q=\frac14n^3 -\frac34n^2+4n-20 = \frac{n^2(n-3)}4+4n-20\]with remainder $100.$ If $n$ is not congruent to $1\pmod4$, then $Q$ is an integer, and
 \[\frac{n^2(n+1)^2}{4} = (n+5)Q + 100 \equiv 17\pmod{n+5},\]so $n+5$ divides $100 - 17 =83$, and $n = 78$. If $n \equiv 1 \pmod4$, then $Q$ is half of an integer, and letting $n = 4k+1$ for some integer $k$ gives
 \[\frac{n^2(n+1)^2}{4} = 2(2k+3)Q + 100 \equiv 17\pmod{n+5}.\]Thus $2k+3$ divides $100-17 = 83$. It follows that $k=40$, and $n = 161$. The requested sum is $161 + 78 = 239$.

 





\textbf{Solution 4}

Using the formula for  $\sum_{k=1}^n  k^3$,
 \[1^3 + 2^3 + 3^3 + ... + n^3 = \frac{n^2(n+1)^2}{4}\]Since $1^3 + 2^3 + 3^3 + ... + n^3$ divided by $n + 5$ has a remainder of $17$,
 \[\frac{n^2(n+1)^2}{4} \equiv 17\pmod {n + 5}\]Using the rules of modular arithmetic,
 \[n^2(n+1)^2 \equiv 68\pmod {n + 5}\]
 \[n^2(n+1)^2 - 68\equiv 0\pmod {n + 5}\]Expanding the left hand side,
 \[n^4 + 2 n^3 + n^2 - 68\equiv 0\pmod {n + 5}\]This means that $n^4 + 2 n^3 + n^2 - 68$ is divisible by ${n + 5}$.

 
 \[(n + 5) | (n^4 + 2 n^3 + n^2 - 68)\]Dividing polynomials,
 \[\frac{n^4 + 2 n^3 + n^2 - 68}{n + 5}\]
 \[= n^3 - 3 n^2 + 16n - 80 + \frac{332}{(n + 5)}\]$(n + 5)$ $|$ $(n^4 + 2 n^3 + n^2 - 68)$ $\iff$ $\frac{332}{(n + 5)}$ $\in$ $\mathbb{Z}$

$\frac{332}{(n + 5)}$ $\in$$\mathbb{Z}$ $\iff$ $(n + 5) = \pm 1, \pm 2, \pm 4, \pm 83, \pm 166, \pm 332$

Note that $n$ $\in$ $\mathbb{N}$ and $n + 5 > 17$ (because the remainder when dividing by $n + 5$ is $17$, so $n + 5$ must be greater than $17$), so all options $\leq 17$ can be eliminated.
 \[(n + 5) = 83, 166, 332\]
 \[n = 78, 161, 327\]Checking all 3 cases, $n = 78$ and $n = 161$ work; $n = 327$ fails.

Therefore, the answer is $78 + 161 = \boxed{239}$.

~ {TSun} ~

 





\textbf{Solution 5 (similar ideas to Solution 1, but faster)}

As before, we note that $(5+a)^3 + (n-a)^3 \equiv (5+a)^3 - (n+5 - (n-a))^3 \equiv 0 \pmod {n+5}.$Thus, we can pair up the terms from $5^3$ to $n^3$ and cancel them. We have to deal with two cases:

 If $n$ is even, then $5^3+6^3 + \cdots + n^3 \equiv 0 \pmod {n+5},$ as there are an even number of terms and they pair and cancel. We thus get $1^2+2^3+3^3+4^3 = 100 \equiv 17 \pmod {n+5},$ or $(n+5) | 83,$ which yields $n=78.$

 If $n$ is odd, then $1^3+2^3+\cdots + n^3 \equiv 1^3+2^3+3^3+4^3+\left( \frac{n+5}{2} \right)^3 \equiv 17 \pmod {n+5}.$Letting $k = \frac{n+5}{2}$ yields $k^2 + 83 \equiv 0 \pmod {2k}.$ However, this means that $83$ is divisible by $k,$ so $k=1,83.$Plugging this back into $n$ yields $n=2(83)-5 = 161$ in the latter case.

 Thus, the sum of all possible $n$ is just $78+161 = \boxed{239}.$

 - ccx09

 \begin{verbatim}
There was a link to alcumus instead of a link to a video solution by OmegaLearn here. I tried to search up his video but I couldn't find it. Please find it and insert it here.
\end{verbatim}




	\item (\textit{AIME II 2020}) Let $P(x) = x^2 - 3x - 7$, and let $Q(x)$ and $R(x)$ be two quadratic polynomials also with the coefficient of $x^2$ equal to $1$. David computes each of the three sums $P + Q$, $P + R$, and $Q + R$ and is surprised to find that each pair of these sums has a common root, and these three common roots are distinct. If $Q(0) = 2$, then $R(0) = \frac{m}{n}$, where $m$ and $n$ are relatively prime positive integers. Find $m + n$.


 
 \textbf{Answer: }071



\textbf{Solution 1}

Let $Q(x) = x^2 + ax + 2$ and $R(x) = x^2 + bx + c$.  We can write the following:
 \[P + Q = 2x^2 + (a - 3)x - 5\]
 \[P + R = 2x^2 + (b - 3)x + (c - 7)\]
 \[Q + R = 2x^2 + (a + b)x + (c + 2)\]Let the common root of $P+Q,P+R$ be $r$; $P+R,Q+R$ be $s$; and $P+Q,Q+R$ be $t$.  We then have that the roots of $P+Q$ are $r,t$, the roots of $P + R$ are $r, s$, and the roots of $Q + R$ are $s,t$.

 By Vieta's, we have:
 \[r + t = \dfrac{3 - a}{2}\tag{1}\]
 \[r + s = \dfrac{3 - b}{2}\tag{2}\]
 \[s + t = \dfrac{-a - b}{2}\tag{3}\]
 \[rt = \dfrac{-5}{2}\tag{4}\]
 \[rs = \dfrac{c - 7}{2}\tag{5}\]
 \[st = \dfrac{c + 2}{2}\tag{6}\]

 Subtracting $(3)$ from $(1)$, we get $r - s = \dfrac{3 + b}{2}$.  Adding this to $(2)$, we get $2r = 3 \implies r = \dfrac{3}{2}$.  This gives us that $t = \dfrac{-5}{3}$ from $(4)$.  Substituting these values into $(5)$ and $(6)$, we get $s = \dfrac{c-7}{3}$ and $s = \dfrac{-3c - 6}{10}$.  Equating these values, we get $\dfrac{c-7}{3} = \dfrac{-3c-6}{10} \implies c = \dfrac{52}{19} = R(0)$.  Thus, our answer is $52 + 19 = \boxed{071}$. ~ TopNotchMath

 





\textbf{Solution 2}

We know that $P(x)=x^2-3x-7$. 

 Since $Q(0)=2$, the constant term in $Q(x)$ is $2$. Let $Q(x)=x^2+ax+2$. 

 Finally, let $R(x)=x^2+bx+c$.

 $P(x)+Q(x)=2x^2+(a-3)x-5$. Let its roots be $p$ and $q$.

 $P(x)+R(x)=2x^2+(b-3)x+(c-7)$ Let its roots be $p$ and $r$.

 $Q(x)+R(x)=2x^2+(a+b)x+(c+2)$. Let its roots be $q$ and $r$.

 By vietas, $p+q=\frac{3-a}{2}, p+r=\frac{3-b}{2}, q+r=\frac{-(a+b)}{2}$

 We could work out the system of equations, but it's pretty easy to see that $p=\frac32, q=-\frac{a}{2}, r=-\frac{b}{2}$.

 $\text{Again, by vietas, }pq=-\frac52\text{, } pr=\frac{c-7}{2}\text{, } qr=\frac{c+2}{2}\text{, } \text{multiplying everything together a}\text{nd taking the sqrt of both sides,}$
 \[(pqr)^2=\left(-\frac52\right)\left(\frac{c-7}{2}\right)\left(\frac{c+2}{2}\right)\]
 \[pqr=\sqrt{\left(-\frac52\right)\left(\frac{c-7}{2}\right)\left(\frac{c+2}{2}\right)}\]$\text{Dividing this }\text{equation by }qr=\frac{c+2}{2}$
 \[\frac{pqr}{qr}=\frac{\sqrt{\left(-\frac52\right)\left(\frac{c-7}{2}\right)\left(\frac{c+2}{2}\right)}}{\frac{c+2}{2}}\]
 \[p = \frac{\sqrt{\left(-\frac52\right)\left(\frac{c-7}{2}\right)}}{\sqrt{\frac{c+2}{2}}}\]$\text{Recall th}\text{at }p=\frac32 \text{ and square both sides}$
 \[\frac94=\frac{\left(-\frac52\right)\left(\frac{c-7}{2}\right)}{\frac{c+2}{2}}\]$\text{Solving gives } c=\frac{52}{19}, \text{ so our answer is }\boxed{071}$

 





\textbf{Solution 3 (Official MAA)}

Let the common root of $P+Q$ and $P+R$ be $p$, the common root of $P+Q$ and $Q+R$ be $q$, and the common root of $Q+R$ and $P+R$ be $r$. Because $p$ and $q$ are both roots of $P+Q$ and $P+Q$ has leading coefficient $2$, it follows that $P(x) + Q(x) = 2(x-p)(x-q).$ Similarly, $P(x) + R(x) = 2(x-p)(x-r)$ and $Q(x) + R(x) = 2(x-q)(x-r)$. Adding these three equations together and dividing by $2$ yields
 \[P(x) + Q(x) + R(x) = (x-p)(x-q) + (x-p)(x-r) + (x-q)(x-r),\]so
 \[P(x) = (P(x) + Q(x) + R(x)) - (Q(x) + R(x))\]
 \[= (x-p)(x-q) + (x-p)(x-r) - (x-q)(x-r)\]
 \[= x^2 - 2px + (pq + pr - qr).\]Similarly,
 \[Q(x) =  x^2 - 2qx + (pq + qr - pr) \text{~ and}\]
 \[R(x) =  x^2 - 2rx + (pr + qr - pq).\]Comparing the $x$ coefficients yields $p = \frac32$, and comparing the constant coefficients yields $-7 = pq + pr - qr = \frac32(q+r) - qr$. The fact that $Q(0) = 2$ implies that $\frac32(q-r) + qr = 2$. Adding these two equations yields $q = -\frac53$, and so substituting back in to solve for $r$ gives $r=-\frac{27}{19}$. Finally,
 \[R(0) = pr + qr - pq = \left(-\frac{27}{19}\right)\left(\frac32-\frac53\right) + \frac52 = \frac{9}{38} + \frac52 = \frac{52}{19}.\]The requested sum is $52 + 19 = 71$. Note that $Q(x) = x^2 + \frac{10}3x + 2$ and $R(x) = x^2 + \frac{54}{19}x + \frac{52}{19}$.

 Another one:

 




	\item (\textit{AIME II 2020}) Let $m$ and $n$ be odd integers greater than $1.$ An $m\times n$ rectangle is made up of unit squares where the squares in the top row are numbered left to right with the integers $1$ through $n$, those in the second row are numbered left to right with the integers $n + 1$ through $2n$, and so on. Square $200$ is in the top row, and square $2000$ is in the bottom row. Find the number of ordered pairs $(m,n)$ of odd integers greater than $1$ with the property that, in the $m\times n$ rectangle, the line through the centers of squares $200$ and $2000$ intersects the interior of square $1099$.


 
 \textbf{Answer: }248



\textbf{Solution 1}

Let us take some cases. Since $m$ and $n$ are odds, and $200$ is in the top row and $2000$ in the bottom, $m$ has to be $3$, $5$, $7$, or $9$. Also, taking a look at the diagram, the slope of the line connecting those centers has to have an absolute value of $< 1$. Therefore, $m < 1800 \mod n < 1800-m$. 

 If $m=3$, $n$ can range from $667$ to $999$. Howe

 




	\item (\textit{AIME II 2020}) Convex pentagon $ABCDE$ has side lengths $AB=5$, $BC=CD=DE=6$, and $EA=7$. Moreover, the pentagon has an inscribed circle (a circle tangent to each side of the pentagon). Find the area of $ABCDE$.


 
 \textbf{Answer: }060



\textbf{Solution 1}

Assume the incircle touches $AB$, $BC$, $CD$, $DE$, $EA$ at $P,Q,R,S,T$ respectively. Then let $PB=x=BQ=RD=SD$, $ET=y=ES=CR=CQ$, $AP=AT=z$. So we have $x+y=6$, $x+z=5$ and $y+z=7$, solve it we have $x=2$, $z=3$, $y=4$. Let the center of the incircle be $I$, by SAS we can proof triangle $BIQ$ is congruent to triangle $DIS$, and triangle $CIR$ is congruent to triangle $SIE$. Then we have $\angle AED=\angle BCD$, $\angle ABC=\angle CDE$. Extend $CD$, cross ray $AB$ at $M$, ray $AE$ at $N$, then by AAS we have triangle $END$ is congruent to triangle $BMC$. Thus $\angle M=\angle N$. Let $EN=MC=a$, then $BM=DN=a+2$. So by law of cosine in triangle $END$ and triangle $ANM$ we can obtain 
 \[\frac{2a+8}{2(a+7)}=\cos N=\frac{a^2+(a+2)^2-36}{2a(a+2)}\], solved it gives us $a=8$, which yield triangle $ANM$ to be a triangle with side length 15, 15, 24, draw a height from $A$ to $NM$ divides it into two triangles with side lengths 9, 12, 15, so the area of triangle $ANM$ is 108. Triangle $END$ is a triangle with side lengths 6, 8, 10, so the area of two of them is 48, so the area of pentagon is $108-48=\boxed{60}$.

 -Fanyuchen20020715

 





\textbf{Solution 2 (Complex Bash)}

Suppose that the circle intersects $\overline{AB}$, $\overline{BC}$, $\overline{CD}$, $\overline{DE}$, and $\overline{EA}$ at $P$, $Q$, $R$, $S$, and $T$ respectively. Then $AT = AP = a$, $BP = BQ = b$, $CQ = CR = c$, $DR = DS = d$, and $ES = ET = e$. So $a + b = 5$, $b + c = 6$, $c + d = 6$, $d + e = 6$, and $e + a = 7$. Then $2a + 2b + 2c + 2d + 2e = 30$, so $a + b + c + d + e= 15$. Then we can solve for each individually. $a = 3$, $b = 2$, $c = 4$, $d = 2$, and $e = 4$. To find the radius, we notice that $4 \arctan(\frac{2}{r}) + 4 \arctan(\frac{4}{r}) + 2 \arctan (\frac{3}{r}) = 360 ^ \circ$, or $2 \arctan(\frac{2}{r}) + 2 \arctan(\frac{4}{r}) + \arctan (\frac{3}{r}) = 180 ^ \circ$. Each of these angles in this could be represented by complex numbers. When two complex numbers are multiplied, their angles add up to create the angle of the resulting complex number. Thus, $(r + 2i)^2 \cdot (r + 4i)^2 \cdot (r + 3i)$ is real. Expanding, we get:

 
 \[(r^2 + 4ir - 4)(r^2 + 8ir -16)(r + 3i)\]

 
 \[(r^4 + 12ir^3 - 52r^2 - 96ir + 64)(r + 3i)\]

 On the last expanding, we only multiply the reals with the imaginaries and vice versa, because we only care that the imaginary component equals 0.

 
 \[15ir^4 - 252ir^2 + 192i = 0\]

 
 \[5r^4 - 84r^2 + 64 = 0\]

 
 \[(5r^2 - 4)(r^2 - 16) = 0\]

 $r$ must equal 4, as r cannot be negative or be approximately equal to 1. Thus, the area of $ABCDE$ is $4 \cdot (a + b + c + d + e) = 4 \cdot 15 = \boxed{60}$

 -nihao4112

 





\textbf{Solution 3 (Guess 1)}

This pentagon is very close to a regular pentagon with side lengths $6$. The area of a regular pentagon with side lengths $s$ is $\frac{5s^2}{4\sqrt{5-2\sqrt{5}}}$. $5-2\sqrt{5}$ is slightly greater than $\frac{1}{2}$ given that $2\sqrt{5}$ is slightly less than $\frac{9}{2}$. $4\sqrt{5-2\sqrt{5}}$ is then slightly greater than $2\sqrt{2}$. We will approximate that to be $2.9$. The area is now roughly $\frac{180}{2.9}$, but because the actual pentagon is not regular, but has the same perimeter of the regular one that we are comparing to we can say that this is an overestimate on the area and turn the $2.9$ into $3$ thus turning the area into $\frac{180}{3}$ which is $60$ and since $60$ is a multiple of the semiperimeter $15$, we can safely say that the answer is most likely $\boxed{60}$.

 





\textbf{Solution 4 (Guess 2)}

Because the AIME answers have to be a whole number it would meant the radius of the circle have to be a whole number, thus by drawing the diagram and experimenting, we can safely say the radius is 4 and the answer is 60

 (Edit: While the guess would be technically correct, the assumption that the radius would have to be a whole number for the ans to be a whole number is wrong)

 By EtherealMidnight

 (Edit: I think that will actually work because the area of $ABCDE$ is equal to the semi-perimeter times the radius. By a simple calculation, we know that the semi-perimeter is an integer so the radius should also be an integer)

 By YBSuburbanTea

 ...the radius could be a fraction with denominator 3, 5, or 15, and the area of the pentagon would still be an integer. - GeometryJake

 





\textbf{Solution 5 (Official MAA 1)}

Let $\omega$ be the inscribed circle, $I$ be its center, and $r$ be its radius. The area of $ABCDE$ is equal to its semiperimeter, $15,$ times $r$, so the problem is reduced to finding $r$. Let $a$ be the length of the tangent segment from $A$ to $\omega$, and analogously define $b$, $c$, $d$, and $e$. Then $a+b=5$, $b+c= c+d=d+e=6$, and $e+a=7$, with a total of $a+b+c+d+e=15$. Hence $a=3$, $b=d=2$, and $c=e=4$. It follows that $\angle B= \angle D$ and $\angle C= \angle E$. Let $Q$ be the point where $\omega$ is tangent to $\overline{CD}$. Then $\angle IAE = \angle IAB =\frac{1}{2}\angle A$. Now we claim that points $A, I, Q$ are collinear, which can be proved if $\angle{AIQ}=\angle{QIA}=180^{\circ}$. The sum of the internal angles in polygons $ABCQI$ and $AIQDE$ are equal, so $\angle IAE + \angle AIQ + \angle IQD + \angle D + \angle E = \angle IAB + \angle B + \angle C + \angle CQI + \angle QIA$, which implies that $\angle AIQ$ must be $180^\circ$. Therefore points $A$, $I$, and $Q$ are collinear.
                    \begin{figure}[H]
                        \centering
                        \includegraphics[width=0.5\textwidth]{Solution/AIME_2020_II_13_Solution_Figure_1.png}
                    \end{figure}
                    Because $\overline{AQ} \perp \overline{CD}$, it follows that
 \[AC^2-AD^2=CQ^2-DQ^2=c^2-d^2=12.\]Another expression for $AC^2-AD^2$ can be found as follows. Note that $\tan \left(\frac{\angle B}{2}\right) = \frac{r}{2}$ and $\tan \left(\frac{\angle E}{2}\right) = \frac{r}{4}$, so
 \[\cos (\angle B) =\frac{1-\tan^2 \left(\frac{\angle B}{2}\right)}{1+\tan^2 \left(\frac{\angle B}{2}\right)} =  \frac{4-r^2}{4+r^2}\]and
 \[\cos (\angle E) = \frac{1-\tan^2 \left(\frac{\angle E}{2}\right)}{1+\tan^2 \left(\frac{\angle E}{2}\right)}= \frac{16-r^2}{16+r^2}.\]Applying the Law of Cosines to $\triangle ABC$ and $\triangle AED$ gives
 \[AC^2=AB^2+BC^2-2\cdot AB\cdot BC\cdot \cos (\angle B) = 5^2+6^2-2 \cdot 5 \cdot 6 \cdot \frac{4-r^2}{4+r^2}\]and
 \[AD^2=AE^2+DE^2-2 \cdot AE \cdot DE \cdot \cos(\angle E) = 7^2+6^2-2 \cdot 7 \cdot 6 \cdot \frac{16-r^2}{16+r^2}.\]Hence

 
 \[12=AC^2- AD^2= 5^2-2\cdot 5 \cdot 6\cdot \frac{4-r^2}{4+r^2} -7^2+2\cdot 7 \cdot 6 \cdot \frac{16-r^2}{16+r^2},\]yielding
 \[2\cdot 7 \cdot 6 \cdot \frac{16-r^2}{16+r^2}- 2\cdot 5 \cdot 6\cdot \frac{4-r^2}{4+r^2}= 36;\]equivalently
 \[7(16-r^2)(4+r^2)-5(4-r^2)(16+r^2) = 3(16+r^2)(4+r^2).\]Substituting $x=r^2$ gives the quadratic equation $5x^2-84x+64=0$, with solutions $\frac{42 - 38}{5}=\frac45$, and $\frac{42 + 38}{5}= 16$. The solution $r^2=\frac45$ corresponds to a five-pointed star, which is not convex. Indeed, if $r<3$, then $\tan \left(\frac{\angle A}{2}\right)$, $\tan \left(\frac{\angle C}{2}\right)$, and $\tan \left(\frac{\angle E}{2}\right)$ are less than $1,$ implying that $\angle A$, $\angle C$, and $\angle E$ are acute, which cannot happen in a convex pentagon. Thus $r^2=16$ and $r=4$. The requested area is $15\cdot4 = \boxed{60}$.

 





\textbf{Solution 6 (Official MAA 2)}

Define $a$, $b$, $c$, $d$, $e$, and $r$ as in Solution 5. Then, as in Solution 5, $a=3$, $b=d=2$, $c=e=4$, $\angle B= \angle D$, and $\angle C= \angle E$. Let $\alpha =\frac{\angle A}{2}$, $\beta = \frac{\angle B}{2}$, and $\gamma=\frac{\angle C}{2}$. It follows that $540^{\circ} = 2\alpha + 4 \beta + 4 \gamma$, so $270^{\circ} = \alpha + 2\beta + 2 \gamma$. Thus
 \[\tan(2\beta + 2 \gamma) = \frac{1}{\tan \alpha},\]$\tan(\beta) = \frac{r}{2}$, $\tan(\gamma) = \frac{r}{4}$, and $\tan(\alpha) = \frac {r}{3}$. By the Tangent Addition Formula,
 \[\tan(\beta +\gamma) = \frac{6r}{8-r^2}\]and
 \[\tan(2\beta + 2\gamma) = \frac{\frac{12r}{8-r^2}}{1-\frac{36r^2}{(8-r^2)^2}} = \frac{12r(8-r^2)}{(8-r^2)^2-36r^2}.\]Therefore
 \[\frac{12r(8-r^2)}{(8-r^2)^2-36r^2} = \frac{3}{r},\]which simplifies to $5r^4 - 84r^2 + 64 = 0$. Then the solution proceeds as in Solution 5.

 





\textbf{Solution 7 (Official MAA 3)}

Define $a$, $b$, $c$, $d$, $e$, and $r$ as in Solution 5. Note that
 \[\arctan\left(\frac{a}{r}\right) + \arctan\left(\frac{b}{r}\right) + \arctan\left(\frac{c}{r}\right) + \arctan\left(\frac{d}{r}\right) + \arctan\left(\frac{e}{r}\right) = 180^{\circ}.\]Hence
 \[\operatorname{Arg}(r + 3i) + 2\cdot \operatorname{Arg}(r + 2i) + 2\cdot  \operatorname{Arg}(r + 4i) = 180^{\circ}.\]Therefore
 \[\operatorname{Im} \big( (r + 3i)(r+2i)^2(r+4i)^2 \big) = 0.\]Simplifying this equation gives the same quadratic equation in $r^2$ as in Solution 5.

 





\textbf{Solution 8 (The same circle)}


                    \begin{figure}[H]
                        \centering
                        \includegraphics[width=0.5\textwidth]{Solution/AIME_2020_II_13_Solution_Figure_2.png}
                    \end{figure}
                    Notation shown on diagram. As in solution 5, we get  $\overline{AQ} \perp  \overline{CD}, AG = 3, GB = 2, CQ = 4$ and so on.

 Let $\overline{AB}$ cross $\overline{CD}$ at $F,  \overline{AE}$ cross $\overline{CD}$ at $F', CF = x.$$FQ = FG \implies FB = x+2.$$\angle BAQ = \angle EAQ \implies DF' = x + 2, EF' = x.$Triangle $\triangle AFF'$ has semiperimeter $s = 2x + 11.$

 The radius of <i>\textbf{incircle}</i> $\omega$ is $r =\sqrt{\frac{s-FF’}{s}}(s-AF) = \sqrt{\frac{3}{2x +11}}(x+4).$

 Triangle $\triangle BCF$ has semiperimeter $s = x + 4.$

 The radius of <i>\textbf{excircle }</i> $\omega$ is $r = \sqrt{\frac{s(s-BF)(s-CF)}{s-BC}} = \sqrt{ \frac{(x+4)\cdot 2 \cdot 4}{x - 2}}.$ 

 It is the same radius, therefore
 \[\sqrt{\frac{3}{2x +11}}(x+4) = \sqrt{\frac{8(x+4)}{x – 2}} \implies \frac {3(x+4)}{2x+11} = \frac {8}{x-2} \implies (x-8)(3x + 14) = 0 \implies x = 8, r = 4.\]

 Then the solution proceeds as in Solution 5.

 





\textbf{Solution 9 (EASIEST)}

Let $O$ be the center of the incircle and $OP$ and $OQ$ are the perpendiculars dropped from point $O$ to sides $BC$ and $CD$ respectively, making $OP = OQ = r$.
After you got $CP$ and $CQ$ by the solutions which are given above (i.e Solution 1) then consider cyclic quadrilateral $OPCQ$.
Now by Pythagoras in TRIANGLE $OPC$
We get $OC^2 = r^2 + 16$, and by Ptolemy you get $4r+4r = OC \cdot PQ$ 
Making $64r^2 = (r^2 + 16)(32)$
 [\textbf{Note: It is unclear how $PQ=\sqrt{32}$ is obtained}; using similar triangles yields $8r=8r$ and there doesn't seem to be any easy way to obtain this otherwise. ~SilverRush] 
leading to $2r^2 = (r^2 + 16)$ implying $r^2 = 16$ implying $r = 4$.
Now since in a tangential polygon Area = $sr$ where s  = semiperimeter of that polygon and r = inradius
Putting those values
\implies AREA = $15r$ = $15(4)$ = $\boxed{60}$

 \textbf{vladimir.shelomovskii@gmail.com, vvsss}

 


 




	\item (\textit{AIME II 2020}) For real number $x$ let $\lfloor x\rfloor$ be the greatest integer less than or equal to $x$, and define $\{x\} = x - \lfloor x \rfloor$ to be the fractional part of $x$. For example, $\{3\} = 0$ and $\{4.56\} = 0.56$. Define $f(x)=x\{x\}$, and let $N$ be the number of real-valued solutions to the equation $f(f(f(x)))=17$ for $0\leq x\leq 2020$. Find the remainder when $N$ is divided by $1000$.


 
 \textbf{Answer: }010



\textbf{Solution 1}

Note that the upper bound for our sum is $2019,$ and not $2020,$ because if it were $2020$ then the function composition cannot equal to $17.$ From there, it's not too hard to see that, by observing the function composition from right to left, $N$ is (note that the summation starts from the right to the left):
 \[\sum_{x=17}^{2019} \sum_{y=x}^{2019} \sum_{z=y}^{2019} 1 .\]One can see an easy combinatorical argument exists which is the official solution, but I will present another solution here for the sake of variety.

 Applying algebraic manipulation and the hockey-stick identity $3$ times gives

 




	\item (\textit{AIME II 2020}) Let $\triangle ABC$ be an acute scalene triangle with circumcircle $\omega$. The tangents to $\omega$ at $B$ and $C$ intersect at $T$. Let $X$ and $Y$ be the projections of $T$ onto lines $AB$ and $AC$, respectively. Suppose $BT = CT = 16$, $BC = 22$, and $TX^2 + TY^2 + XY^2 = 1143$. Find $XY^2$.


 
 \textbf{Answer: }717



\textbf{Solution 1}

Let $O$ be the circumcenter of $\triangle ABC$; say $OT$ intersects $BC$ at $M$; draw segments $XM$, and $YM$.  We have $MT=3\sqrt{15}$. 

 
                    \begin{figure}[H]
                        \centering
                        \includegraphics[width=0.5\textwidth]{Solution/AIME_2020_II_15_Solution_Figure_1.png}
                    \end{figure}
                    Since $\angle A=\angle CBT=\angle BCT$, we have $\cos A=\frac{11}{16}$. Notice that $AXTY$ is cyclic, so $\angle XTY=180^{\circ}-A$, so $\cos XTY=-\cos A$, and the cosine law in $\triangle TXY$ gives 
 \[1143-2XY^2=-\frac{11}{8}\cdot XT\cdot YT.\] 

 Since $\triangle BMT \cong \triangle CMT$, we have $TM\perp BC$, and therefore quadrilaterals $BXTM$ and $CYTM$ are cyclic. Let $P$ (resp. $Q$) be the midpoint of $BT$ (resp. $CT$). So $P$ (resp. $Q$) is the center of $(BXTM)$ (resp. $CYTM$). Then $\theta=\angle ABC=\angle MTX$ and $\phi=\angle ACB=\angle YTM$. So $\angle XPM=2\theta$, so
 \[\frac{\frac{XM}{2}}{XP}=\sin \theta,\]which yields $XM=2XP\sin \theta=BT(=CT)\sin \theta=TY$. Similarly we have $YM=XT$. 

 Ptolemy's theorem in $BXTM$ gives 
 \[16TY=11TX+3\sqrt{15}BX,\] while Pythagoras' theorem gives $BX^2+XT^2=16^2$. Similarly, Ptolemy's theorem in $YTMC$ gives
 \[16TX=11TY+3\sqrt{15}CY\] while Pythagoras' theorem in $\triangle CYT$ gives $CY^2+YT^2=16^2$. Solve this for $XT$ and $TY$ and substitute into the equation about $\cos XTY$ to obtain the result $XY^2=\boxed{717}$.

 (Notice that $MXTY$ is a parallelogram, which is an important theorem in Olympiad, and there are some other ways of computation under this observation.)

 -Fanyuchen20020715

 





\textbf{Solution 2 (Official MAA)}

Let $M$ denote the midpoint of $\overline{BC}$. The critical claim is that $M$ is the orthocenter of $\triangle AXY$, which has the circle with diameter $\overline{AT}$ as its circumcircle. To see this, note that because $\angle BXT = \angle BMT = 90^\circ$, the quadrilateral $MBXT$ is cyclic, it follows that
 \[\angle MXA = \angle MXB = \angle MTB = 90^\circ - \angle TBM = 90^\circ - \angle A,\] implying that $\overline{MX} \perp \overline{AC}$. Similarly, $\overline{MY} \perp \overline{AB}$. In particular, $MXTY$ is a parallelogram.
                    \begin{figure}[H]
                        \centering
                        \includegraphics[width=0.5\textwidth]{Solution/AIME_2020_II_15_Solution_Figure_2.png}
                    \end{figure}
                    Hence, by the Parallelogram Law,
 \[TM^2 + XY^2 = 2(TX^2 + TY^2) = 2(1143-XY^2).\] But $TM^2 = TB^2 - BM^2 = 16^2 - 11^2 = 135$. Therefore 
 \[XY^2 = \frac13(2 \cdot 1143-135) = \boxed{717}.\]

 





\textbf{Solution 3 (Law of Cosines)}

Let $H$ be the orthocenter of $\triangle AXY$. 

 \textbf{Lemma 1:} $H$ is the midpoint of $BC$.

 \textbf{Proof:} Let $H'$ be the midpoint of $BC$, and observe that $XBH'T$ and $TH'CY$ are cyclical. Define $H'Y \cap BA=E$ and $H'X \cap AC=F$, then note that:
 \[\angle H'BT=\angle H'CT=\angle H'XT=\angle H'YT=\angle A.\]That implies that $\angle H'XB=\angle H'YC=90^\circ-\angle A$, $\angle CH'Y=\angle EH'B=90^\circ-\angle B$, and $\angle BH'Y=\angle FH'C=90^\circ-\angle C$. Thus $YH'\perp AX$ and $XH' \perp AY$; $H'$ is indeed the same as $H$, and we have proved lemma 1.

 Since $AXTY$ is cyclical, $\angle XTY=\angle XHY$ and this implies that $XHYT$ is a paralelogram.By the Law of Cosines:
 \[XY^2=XT^2+TY^2+2(XT)(TY)\cdot \cos(\angle A)\]
 \[XY^2=XH^2+HY^2+2(XH)(HY) \cdot \cos(\angle A)\]
 \[HT^2=HX^2+XT^2-2(HX)(XT) \cdot \cos(\angle A)\]
 \[HT^2=HY^2+YT^2-2(HY)(YT) \cdot \cos(\angle A).\]We add all these equations to get:
 \[HT^2+XY^2=2(XT^2+TY^2) \qquad (1).\]We have that $BH=HC=11$ and $BT=TC=16$ using our midpoints. Note that $HT \perp BC$, so by the Pythagorean Theorem, it follows that $HT^2=135$. We were also given that $XT^2+TY^2=1143-XY^2$, which we multiply by $2$ to use equation $(1)$. 
 \[2(XT^2+TY^2)=2286-2 \cdot XY^2\] Since $2(XT^2+TY^2)=2(HT^2+TY^2)=HT^2+XY^2$, we have
 \[135+XY^2=2286-2 \cdot XY^2\] 
 \[3 \cdot XY^2=2151.\]Therefore, $XY^2=\boxed{717}$. ~ MathLuis

 





\textbf{Solution 4 (Similarity and median)}


                    \begin{figure}[H]
                        \centering
                        \includegraphics[width=0.5\textwidth]{Solution/AIME_2020_II_15_Solution_Figure_3.png}
                    \end{figure}
                    Using the <i>\textbf{Claim}</i> (below) we get $\triangle ABC \sim \triangle XTM \sim \triangle YMT.$

 Corresponding sides of similar $\triangle XTM \sim \triangle YMT$ is $MT,$ so

 $\triangle XTM = \triangle YMT \implies MY = XT, MX = TY \implies XMYT$ – parallelogram.

 
 \[4 TD^2 = MT^2 = \sqrt{BT^2 - BM^2}  =\sqrt{153}.\]The formula for median $DT$ of triangle $XYT$ is
 \[2 DT^2 = XT^2 + TY^2 – \frac{XY^2}{2},\]
 \[3 \cdot XY^2 = 2XT^2 + 2TY^2 + 2XY^2 –  4 DT^2,\]
 \[3 \cdot XY^2  = 2 \cdot 1143-153 = 2151 \implies XY^2 = \boxed{717}.\]

 


 
                    \begin{figure}[H]
                        \centering
                        \includegraphics[width=0.5\textwidth]{Solution/AIME_2020_II_15_Solution_Figure_4.png}
                    \end{figure}
                    <i>\textbf{Claim}</i>

 Let $\triangle ABC$ be an acute scalene triangle with circumcircle $\omega$. The tangents to $\omega$ at $B$ and $C$ intersect at $T$. Let $X$ be the projections of $T$ onto line $AB$. Let M be midpoint BC. Then triangle ABC is similar to triangle XTM.

 <i>\textbf{Proof}</i>

 $\angle BXT = \angle BMT = 90^o \implies$ the quadrilateral $MBXT$ is cyclic.

 $BM \perp MT, TX \perp AB \implies \angle MTX = \angle MBA.$

 $\angle CBT = \angle BAC = \frac {\overset{\Large\frown} {BC}}{ 2} \implies \triangle ABC \sim \triangle XTM.$

 \textbf{vladimir.shelomovskii@gmail.com, vvsss}

 





\textbf{Solution 5 (Symmedian and Stewarts)}

Let $M$ be the midpoint of $BC$. Note that $\angle XYT = \angle XAT = \angle MAC$ because $AT$ is a symmedian. Similarly, $\angle TXY = \angle MAB$. $XT = 16\sin{C}$ and $YT = 16\sin{C}$. By law of sines on $XYT$, $\frac{16\sin{C}}{\sin{XYT}} = \frac{XY}{\sin{A}}$. However by law of sines on $MAC$, $\frac{\sin{C}}{\sin{MAC}} = \frac{AM}{11}$. Combining these two yields, $\frac{16}{11} AM = \frac{XY}{\sin{A}}$. Since $\sin{A} = \frac{\sqrt{135}}{16}$, we have $XY^2 = \frac{135}{121} AM^2$.

 Letting $AB = x$, $AC = y$, and $AM = d$, we have $x^2 + y^2 = 2d^2 + 242$ by Stewarts. Since $x = 2R\sin{C}$, and $y = 2R\sin{B}$ by extended law of sines, we can write $\sin^2{B} + \sin^2{C} = \frac{2d^2 + 242}{4R^2}$. By law of cosines on $ABC$, $x^2 + y^2 - 2xy(\frac{11}{16}) = 484$, $\frac{11}{8} xy + 484 = 2d^2 + 242$, $xy = \frac{8}{11} (2d^2 - 242)$. Then similarly as before we can write $\sin{B}\sin{C} = \frac{\frac{8}{11} (2d^2 - 242)}{4R^2}$.

 By law of cosines on $XYT$ and using $XT^2 + XY^2 + YT^2 = 1143$, we have $512(\sin^2{B} + \sin^2{C}) - 1143 = -512\sin{B}\sin{C}(\frac{11}{16})$. Substituting our previous values here and using $R = \frac{176}{\sqrt{135}}$ yields a value for $d^2$, and multiplying by $\frac{135}{121}$ gives $\boxed{717}$.

 




\end{enumerate}